\documentclass[10pt,a4paper]{amsart}
\usepackage[margin=19mm]{geometry}
\usepackage{amsmath,amssymb,mathtools}
\usepackage[scr=rsfs,cal=euler]{mathalfa}
\usepackage{xfrac}
\usepackage[unicode,hidelinks,bookmarks=false]{hyperref}
\newcommand{\ie}{i.e.\ }
\newcommand{\cf}{cf.\ }
\newcommand{\resp}{respectively\ }
\newcommand{\Subsection}{\subsection}
\providecommand{\nobreakdash}{\nobreak\hbox{-}}
\setlength{\emergencystretch}{2em}
\multlinegap0pt
\usepackage{amssymb}
\usepackage{bm}
\newtheorem{theorem}{Theorem}
\def\T{{ \mathrm{\scriptscriptstyle T} }}
\def\pr{{\operatorname{pr}}}
\title{Gaussian comparison above the median}
\author{Colin B. Fogarty}
\date{Source: arXiv:2607.06874v2 (CC BY 4.0)}
\begin{document}
\maketitle

\noindent\emph{Source and licence.} Gaussian comparison above the median. Version arXiv:2607.06874v2 (CC BY 4.0), distributed by arXiv under the Creative Commons Attribution 4.0 International licence, http://creativecommons.org/licenses/by/4.0/, as recorded in the arXiv licence metadata for this exact version and shown on the abstract page https://arxiv.org/abs/2607.06874v2. Full licence notice: \texttt{licenses/arxiv-2607.06874-CC-BY-4.0.txt}.\par\smallskip

\noindent\textit{Editorial scope.} Counted unit: the article's theorem thm:conv (source0.tex lines 156--161). The article's complete original proof of it is delivered with it (source0.tex lines 262--290, one \texttt{proof} environment of 29 lines, which the article places away from the statement). No mathematics is written, repaired or re-derived: the statement and the proof are the article's own lines, in the article's own notation, and no retained line is substituted or re-flowed.  No further source-local context is needed: every cross-reference of every retained line resolves inside the two delivered spans.  Nothing was omitted from any retained span, so the package carries no omission record.  No retained line contains a citation, so the wrapper carries no bibliography and no bibliography entry.  No retained line contains a figure environment, an image-inclusion command or drawing code, so no image is delivered and the package declares no asset dependency.  Editorial formatting only, in addition: an AMS \texttt{amsart} preamble that restates the article's own declarations and macros used by the retained lines, namely the environments \texttt{theorem} on separate counters, together with the standard packages those lines need.  The retained environments print in this document as Theorem 1 in order of appearance, whereas the article prints the same items with the numbering of its own full document.  The complete native source of the article as distributed in the arXiv e-print for this exact version is bundled unchanged as \texttt{sources/204667/source0.tex}.
\begin{theorem}\label{thm:conv}
Let $X\in \mathbb{R}^d$ be a mean-zero Gaussian random vector with covariance $\Sigma_X\succ 0$, and let $W \in \mathbb{R}^d$ be a random vector independent of $X$ satisfying $W \overset{d}{=}-W$. For every closed convex set $\mathcal{K}$,
\begin{align*}
\pr(X+W \in \mathcal{K})\geq \frac{1}{2} \Rightarrow \pr(X \in \mathcal{K})\geq \pr(X+W\in \mathcal{K}).
\end{align*}
\end{theorem}

\begin{proof}
We break the proof into two cases, depending on whether or not $\Sigma_Y$ is positive definite.

\medskip
\noindent\textit{Case 1: $\Sigma_Y$ is positive definite.}

Let $W \sim \mathcal{N}_d(0, \Sigma_Y-\Sigma_X)$ with $W$ independent of $X$. Observe that $Y \overset{d}{=}X+W$ and that $W\overset{d}{=} -W$.
Suppose first that $\text{rank}(\Sigma_X) = d$. Then, Theorem \ref{thm:conv} implies that if $\pr(Y\in \mathcal{K}) = \pr(X+W \in \mathcal{K})\geq 1/2$, 
\begin{align*} \pr(X\in \mathcal{K})&\geq \pr(X+W\in \mathcal{K})\\ &= \pr(Y\in \mathcal{K}).\end{align*}

Now suppose $\text{rank}(\Sigma_X) < d$. For $t \in [0,1]$, define 
\begin{align*} Z_t &= X + t^{1/2}W.
\end{align*} For all $t\in(0,1]$, $Z_t \sim \mathcal{N}_d(0,\Sigma_t)$ with 
\begin{align*} \Sigma_t &= \Sigma_X+t(\Sigma_Y-\Sigma_X) \succ 0;\\ \Sigma_Y-\Sigma_t &= (1-t)(\Sigma_Y-\Sigma_X) \succeq 0.\end{align*} For each $t\in (0,1]$, let $W_t$ be independent of $Z_t$ with 
\begin{align*}W_t \sim \mathcal{N}_d(0, \Sigma_Y-\Sigma_t),
\end{align*}so that $Z_t+W_t \overset{d}{=} Y$ and $W_t \overset{d}{=} -W_t$. The preceding full-rank argument then establishes that for $t \in (0,1]$, whenever $\pr(Y\in \mathcal{K})\geq 1/2$,\begin{align*}\pr(Z_t \in \mathcal{K}) \geq \pr(Y\in\mathcal{K}).\end{align*}Observe that $||X - Z_t||_2^2 = t||W||_2^2$. As $E||W||^2_2 < \infty$, we have that $Z_t$ converges in $L^2$, and hence in distribution, to $X$ as $t\downarrow 0$. As $\mathcal{K}$ is closed, the Portmanteau theorem implies 
\begin{align*}\operatorname{pr}(X\in \mathcal{K}) \geq \limsup_{t\downarrow 0}\operatorname{pr}(Z_t \in \mathcal{K})\geq \operatorname{pr}(Y\in \mathcal{K})\end{align*} as desired.

\medskip
\noindent\textit{Case 2: $\Sigma_Y$ is singular.}

Let $r := \text{rank}(\Sigma_Y)$. If $r=0$ the result is immediate, so suppose $r\geq 1$. Observe that as $0\preceq\Sigma_X\preceq\Sigma_Y$, any $z$ in the null space of $\Sigma_Y$ must be in the null space of $\Sigma_X$, as
\begin{align*} 0\leq z^\T\Sigma_X z\leq z^\T\Sigma_Y z = 0.\end{align*} As $\Sigma_X\succeq 0$ , this implies $\Sigma_Xz = 0$. Thus, by symmetry of $\Sigma_X$ and $\Sigma_Y$, 
\begin{align*}\text{colspace}(\Sigma_X)\subseteq \text{colspace}(\Sigma_Y).\end{align*} Eigendecompose $\Sigma_Y$ as $\tilde{Q}\tilde{\Lambda} \tilde{Q}^\T$. $\tilde{\Lambda}$ is diagonal, with $r$ positive entries and $d-r$ entries equal to zero. Partition $\tilde{Q}=[U,U_0]$ for the eigenvectors corresponding to the $r$ positive and $d-r$ null eigenvalues respectively. Define 
\begin{align*} \tilde{X} = U^\T X \in \mathbb{R}^r;\\\tilde{Y} = U^\T Y \in \mathbb{R}^r,\end{align*} and observe that 
\begin{align*} \Sigma_{\tilde{Y}}\succ 0;\;\;\;\Sigma_{\tilde{Y}} - \Sigma_{\tilde{X}} = U^\T(\Sigma_Y-\Sigma_X)U \succeq 0.\end{align*}Define $\tilde{\mathcal{K}} = \{q\in\mathbb{R}^r: U q \in \mathcal{K}\}$.  Then, since $X, Y \in \text{colspace}(\Sigma_Y)$ a.s., 
\begin{align*}\operatorname{pr}(\tilde{X}\in \tilde{\mathcal{K}}) &= \operatorname{pr}(X\in \mathcal{K});\\ \operatorname{pr}(\tilde{Y}\in \tilde{\mathcal{K}}) &= \operatorname{pr}(Y\in \mathcal{K})\geq 1/2.\end{align*} 
Since $\mathcal{K}$ is closed and convex and $q\mapsto Uq$ is linear, $\tilde{\mathcal{K}}$ is also closed and convex. The proof then follows by applying the Case 1 argument to $\tilde{X}$, $\tilde{Y}$, and $\tilde{\mathcal{K}}$. 
\end{proof}

\end{document}
